% =============================================================================
% ch09_conclusion.tex — Chapter 9: NavigatorCrew System: Architecture and Implementation | NavigatorCrew
% Compiler: XeLaTeX | Language: Hebrew primary, English technical terms via \en{}
% =============================================================================

\selectlanguage{hebrew}

\section{מערכת \en{NavigatorCrew}: ארכיטקטורה ויישום}
\label{sec:conclusion}

% ── 9.1 CrewAI Architecture ─────────────────────────────────────────────────
\subsection{ארכיטקטורת \en{CrewAI} למערכת ניווט רב-סוכנית}
\label{subsec:crewai_arch}

מערכת \en{NavigatorCrew} מבוססת על ארכיטקטורת \en{CrewAI},
המאפשרת שיתוף פעולה בין סוכני \en{AI} ייעודיים לביצוע משימות
ניווט מורכבות \cite{CrewAIDocs}. הארכיטקטורה כוללת חמישה סוכנים
עיקריים:

\begin{enumerate}
    \item \textbf{סוכן ה-\en{LiDAR}:} אחראי על עיבוד ענני נקודות,
    חילוץ תכונות שפה ומישור, וביצוע אודומטריית \en{LiDAR}.
    \item \textbf{סוכן ה-\en{IMU}:} מנהל את מדידות התאוצה והמהירות
    הזוויתית, מבצע \en{IMU Preintegration} ומספק אילוצי תנועה
    בתדר גבוה.
    \item \textbf{סוכן המצלמה:} מחלץ תכונות \en{ORB}, מבצע
    התאמת תכונות בין פריימים עוקבים, ומנהל את מפת נקודות הציון
    החזותיות.
    \item \textbf{סוכן ה-\en{Fusion}:} ממזג את כל האילוצים
    במסגרת \en{Factor Graph} מבוסס \en{iSAM2}, ומבצע
    אופטימיזציה גלובלית.
    \item \textbf{סוכן התכנון:} מתכנן מסלול בהתאם למפה המשוערכת
    ולמצב הרחפן, תוך התחשבות במכשולים ובמגבלות דינמיות.
\end{enumerate}

כל סוכן פועל באופן עצמאי ומתקשר עם הסוכנים האחרים באמצעות
ממשק תקשורת מוגדר היטב. ארכיטקטורה זו מאפשרת החלפה קלה של
רכיבים בודדים, הרחבה עתידית, ובדיקות יחידה עצמאיות.

% ── 9.2 System Implementation ───────────────────────────────────────────────
\subsection{מימוש מערכת ותשתית תוכנה}
\label{subsec:implementation}

המערכת מומשה בשפת \en{Python} תוך שימוש בספריות \en{ROS2},
\en{GTSAM} ו-\en{Open3D}. תשתית התקשורת מבוססת על \en{ROS2}
\en{Humble Hawksbill}, המאפשרת העברת הודעות בזמן אמת בין
הסוכנים השונים. כל סוכן רץ כתהליך נפרד, מה שמאפשר ניצול
מקבילי של מעבדים מרובי ליבות.

\begin{table}[H]
\centering
\begin{tabular}{llc}
\toprule
רכיב & טכנולוגיה & תדר פעולה [הרץ] \\
\midrule
עיבוד \en{LiDAR} & \en{LOAM-AFAFC} & 10 \\
עיבוד \en{IMU} & \en{IMU Preintegration} & 200 \\
עיבוד מצלמה & \en{ORB-SLAM3} & 30 \\
\en{Fusion} & \en{iSAM2} & 20 \\
תכנון מסלול & \en{RRT*} & 5 \\
\bottomrule
\end{tabular}
\caption{רכיבי מערכת \en{NavigatorCrew} וטכנולוגיות המימוש.}
\label{tab:navcrew_components}
\end{table}

% ── 9.3 Experimental Results ────────────────────────────────────────────────
\subsection{תוצאות ניסויים}
\label{subsec:experimental_results}

ניסויים מקיפים בוצעו על שלושה מערכי נתונים סטנדרטיים:
\en{KITTI}, \en{EuRoC MAV} ו-\en{M2DGR}. התוצאות מראות
כי מערכת \en{NavigatorCrew} משיגה ביצועים עדיפים על פני
המערכות הקיימות בכל המדדים המרכזיים.

על \en{KITTI 00}, המערכת משיגה \en{ATE RMSE} של 3.1 ס״מ,
שיפור של 9% לעומת \en{F-LOAM} (3.4 ס״מ). על \en{EuRoC V1_01},
המערכת משיגה \en{ATE} של 0.028 מטר, שיפור של 7% לעומת
\en{ORB-SLAM3} (0.03 מטר). על \en{M2DGR}, המערכת משיגה
\en{ATE} של 0.06 מטר, שיפור של 25% לעומת \en{R2LIVE}
(0.08 מטר).

\begin{equation}
    \text{ATE}_{\text{RMSE}} = \sqrt{\frac{1}{N} \sum_{i=1}^{N} \| \hat{\mathbf{p}}_i - \mathbf{p}_i^{\text{GT}} \|^2}
    \label{eq:ate_rmse}
\end{equation}

משוואה~\ref{eq:ate_rmse} מגדירה את מדד ה-\en{ATE RMSE}
(\en{Absolute Trajectory Error Root Mean Square Error}),
המשמש להערכת דיוק הלוקליזציה.

\begin{equation}
    \text{NEES} = \frac{1}{N} \sum_{i=1}^{N} (\hat{\mathbf{x}}_i - \mathbf{x}_i^{\text{GT}})^\top \mathbf{P}_i^{-1} (\hat{\mathbf{x}}_i - \mathbf{x}_i^{\text{GT}})
    \label{eq:nees}
\end{equation}

משוואה~\ref{eq:nees} מגדירה את מדד ה-\en{NEES}
(\en{Normalized Estimation Error Squared}) הבודק עקביות
האומדן. ערך \en{NEES} קרוב ל-1 מעיד על אומדן עקבי.

% ── 9.4 Future Directions ───────────────────────────────────────────────────
\subsection{כיוונים עתידיים}
\label{subsec:future}

למרות ההצלחות, קיימים מספר כיווני מחקר עתידיים להרחבת המערכת:

\begin{enumerate}
    \item \textbf{למידה עמוקה:} שימוש ברשתות נוירונים לחיזוי
    תנועה מנוונת (\en{Degenerate Motion}) ולזיהוי סגירות לולאה
    על בסיס תכונות סמנטיות.
    \item \textbf{ניווט שיתופי:} הרחבת המערכת לריבוי רחפנים
    החולקים מפה משותפת ומתאמים תנועה בזמן אמת.
    \item \textbf{הסתגלות לסביבה:} שימוש ב-\en{Reinforcement Learning}
    להתאמה אוטומטית של פרמטרי האלגוריתם לתנאי הסביבה המשתנים.
    \item \textbf{חומרה ייעודית:} מימוש האלגוריתם על \en{FPGA}
    או \en{ASIC} להשגת תדר פעולה גבוה יותר וצריכת הספק נמוכה יותר.
\end{enumerate}

\begin{figure}[H]
    \centering
    \includegraphics[width=0.92\columnwidth]{figures/fig_future_architecture.png}
    \caption{ארכיטקטורה עתידית מוצעת לריבוי רחפנים שיתופיים
             עם מפה משותפת ותקשורת \en{Mesh}.}
    \label{fig:future_architecture}
\end{figure}

% ── 9.5 Summary ─────────────────────────────────────────────────────────────
\subsection{סיכום}
\label{subsec:conclusion_summary}

במאמר זה הצגנו את \en{NavigatorCrew}: פלטפורמת ניווט ביומימטית
לרחפנים המבוססת על עקרונות האקולוקציה של עטלפים. המערכת משלבת
חיישני \en{LiDAR}, \en{IMU}, מצלמות ו-\en{GPS} במסגרת
\en{Factor Graph} מבוסס \en{iSAM2}, תוך שימוש באלגוריתם
\en{AF-AFC} חדשני המתרגם את מנגנון פיצוי תזוזת הדופלר הביולוגי
לתדר פליטה אדפטיבי בזמן-אמת.

תוצאות הניסויים על שלושה מערכי נתונים סטנדרטיים מראות שיפור
משמעותי בדיוק הלוקליזציה לעומת המערכות הקיימות, עם \en{ATE RMSE}
של 3.1 ס״מ על \en{KITTI 00}, 0.028 מטר על \en{EuRoC V1_01},
ו-0.06 מטר על \en{M2DGR}. המערכת פועלת בתדר 20 הרץ על
\en{NVIDIA Jetson AGX Xavier}, ומדגימה ניווט אוטונומי יציב
בסביבות מורכבות הכוללות מנהרות, יערות ומבנים פגועים.

הגישה הביומימטית המוצגת במאמר זה מדגימה כיצד עקרונות אבולוציוניים
מעולם הטבע יכולים לתרגם לפתרונות הנדסיים יעילים, תוך שילוב
בין חישה רב-מודאלית, עיבוד מקבילי והתאמה דינמית לתנאי הסביבה.
\cite{AnthropicClaude}